\documentclass[12pt]{article}
\input{iati_structure.tex}
\renewcommand{\bottomtitlespace}{3cm}

\usepackage[english]{babel}

\begin{document}

\renewcommand{\headerLang}{Ελληνικά}
\renewcommand{\problemName}{AB13. Vision}

\renewcommand{\headerLeft}{IATI Μέρα 1, Senior κατηγορία}
\renewcommand{\headerRightFirst}{XVII INTERNATIONAL ADVANCED TOURNAMENT IN INFORMATICS}
\renewcommand{\headerRightSecond}{ONLINE 2026}

\renewcommand{\tl}{$5$ s}
\renewcommand{\ml}{$1024$ MB}
\problem{Πρόβλημα \problemName}
\addauthor{Author: Emil Indzhev}{2026}{04}{16}

Το διαστημόπλοιο USS Enterprise (με τον Sashka ως νέο κυβερνήτη του) ταξιδεύει σε όλο το σύμπαν, το οποίο αποτελεί ένα άπειρο δισδιάστατο πλέγμα τομέων. Σε κάθε τομέα $(X, Y)$\footnote{Το $X$ είναι ο αριθμός γραμμής και αυξάνεται προς τα κάτω, ενώ το $Y$ είναι ο αριθμός στήλης και αυξάνεται προς τα δεξιά.} υπάρχει ένας φάρος με ισχύ όρασης $V_{X, Y}$. Όταν το σκάφος βρίσκεται σε αυτόν τον τομέα, χρησιμοποιεί τον φάρο για να σαρώσει όλους τους τομείς $(X', Y')$ με $X - V_{X, Y} \leq X' \leq X + V_{X, Y}$ και  $Y - V_{X, Y} \leq Y' \leq Y + V_{X, Y}$. Έτσι, οι ορατοί τομείς σχηματίζουν ένα τετράγωνο πλευράς $2V_{X, Y} + 1$. Για κάθε τέτοιο τομέα, η μόνη πληροφορία που είναι ορατή είναι η ισχύς όρασης του φάρου σε αυτόν, δηλαδή $V_{X', Y'}$ (επειδή οι φάροι λειτουργούν ταυτόχρονα ως πομποί και ως δέκτες). Μετά τη σάρωση, το σκάφος τηλεμεταφέρεται σε έναν από αυτούς τους ορατούς τομείς, καθώς οι φάροι λειτουργούν επίσης ως συσκευές τηλεμεταφοράς.

Δυστυχώς, οι φάροι έχουν ένα κρίσιμο ελάττωμα: μπορεί να αντιστρέψουν τον άξονα $X$ της σάρωσης, ή τον άξονα $Y$, ή και τους δύο (δηλαδή υπάρχουν $4$ δυνατοί προσανατολισμοί της σάρωσης). Παρατηρήστε ότι το σκάφος χρησιμοποιεί αυτήν την, ενδεχομένως αντεστραμμένη, σάρωση για να αποφασίσει σε ποιον τομέα θα τηλεμεταφερθεί, αλλά και η ίδια η τηλεμεταφορά γίνεται με τον ίδιο προσανατολισμό, π.χ. αν ο άξονας $Y$ είναι αντεστραμμένος και το σκάφος επιλέξει να τηλεμεταφερθεί προς τα αριστερά στη σάρωση (μειώνοντας το $Y$), τότε στην πραγματικότητα θα τηλεμεταφερθεί προς τα δεξιά. Αυτό σημαίνει ότι το σκάφος πράγματι θα τηλεμεταφερθεί στον τομέα που επέλεξε μέσα στη σάρωση.

Επιπλέον, το σκάφος δεν έχει μνήμη, οπότε η απόφασή του για το πού θα τηλεμεταφερθεί πρέπει να εξαρτάται μόνο από τη σάρωση που παράγεται από τον φάρο. Ακόμη, η στρατηγική του πρέπει να είναι ντετερμινιστική -- αν του δοθεί η ίδια σάρωση, πρέπει να επιλέγει πάντα τον ίδιο τομέα μέσα σε αυτή.

Το σκάφος ξεκινά από άγνωστες συντεταγμένες και πρέπει πάντοτε (ανεξάρτητα από το πού ξεκινά και ανεξάρτητα από το ποιοι προσανατολισμοί σάρωσης εμφανίζονται) να μπορεί να ταξιδεύει προς την κατεύθυνση αύξησης των $X$ και $Y$ όσο μακριά θέλει (μπορούμε να πούμε ότι απαιτούμε τελικά να φτάσει σε κάποιον τομέα $(X, Y)$ με $X, Y \geq 10^{100}$). Ωστόσο, για να είναι αυτό δυνατό, είναι καθοριστικής σημασίας τόσο η στρατηγική μετακίνησης του σκάφους όσο και η διάταξη των τιμών όρασης σε όλο το σύμπαν.

Εδώ αναλαμβάνετε εσείς -- πρέπει να σχεδιάσετε ένα ορθογώνιο πρότυπο (μοτίβο) $P$ διαστάσεων $N$ επί $M$ με τιμές όρασης, το οποίο θα επαναλαμβάνεται άπειρα σε όλο το σύμπαν: $V_{X, Y} = P_{X \bmod N, Y \bmod M}$. Πρέπει επίσης να σχεδιάσετε τη στρατηγική μετακίνησης του σκάφους, δηλαδή μια συνάρτηση που δέχεται ως είσοδο τη σάρωση που παράγεται από τον φάρο και επιλέγει σε ποιον από τους ορατούς τομείς θα τηλεμεταφερθεί το σκάφος στη συνέχεια.

Επειδή οι φάροι με μεγαλύτερη ισχύ όρασης είναι ακριβοί, στόχος σας είναι να δώσετε μια έγκυρη λύση, προσπαθώντας παράλληλα να ελαχιστοποιήσετε τη μέση ισχύ όρασης στο πρότυπό σας, δηλαδή τη μέση τιμή των στοιχείων του $P$.

Επειδή το πρόβλημα αυτό ίσως σας φαίνεται αρκετά δύσκολο, αποφασίζετε να ξεκινήσετε με ένα πρόβλημα εξάσκησης: το ίδιο πρόβλημα αλλά σε μία διάσταση. Σε αυτήν την απλούστερη εκδοχή, αγνοούμε τη διάσταση $X$ και χρησιμοποιούμε μόνο τη διάσταση $Y$. Έτσι, το πρότυπό σας είναι απλώς μια ακολουθία μήκους $M$, η οποία επαναλαμβάνεται ως εξής: $V_Y = P_{Y \bmod M}$. Οι σαρώσεις που παράγονται από τους φάρους είναι επίσης ακολουθίες μήκους $2V_Y + 1$ (που περιέχουν τους τομείς $Y'$ για τους οποίους ισχύει  $Y - V_Y \leq Y' \leq Y + V_Y$). Εδώ υπάρχουν μόνο $2$ δυνατοί προσανατολισμοί της σάρωσης (κανονικός ή αντεστραμμένος) και το σκάφος πρέπει να μπορεί να φτάσει σε θέση με $Y \geq 10^{100}$.

Οι λύσεις σας για τη μονοδιάστατη και τη δισδιάστατη εκδοχή θα βαθμολογηθούν ανεξάρτητα, και μπορείτε να πάρετε βαθμούς για τη μία ακόμη και χωρίς να έχετε έγκυρη λύση για την άλλη.

\subsection{Λεπτομέρειες υλοποίησης}

Πρέπει να υλοποιήσετε $4$ συναρτήσεις: \verb|getVisionPattern1d|, \verb|getMove1d|, \verb|getVisionPattern2d| και \verb|getMove2d|. Οι πρώτες $2$ χρησιμοποιούνται για τη μονοδιάστατη εκδοχή και οι άλλες $2$ για τη δισδιάστατη. Σε κάθε τεστ ελέγχου θα κληθεί μόνο ένα από αυτά τα δύο ζεύγη συναρτήσεων, όμως και οι $4$ πρέπει να έχουν υλοποιηθεί (ακόμη κι αν οι υλοποιήσεις του ενός ζεύγους είναι κενές) ώστε η λύση να θεωρείται έγκυρη (να μεταγλωττίζεται).

\begin{lstlisting}
 std::vector<int> getVisionPattern1d()
\end{lstlisting}

Η συνάρτηση αυτή θα κληθεί μία φορά, αν το τεστ ελέγχου αφορά τη μονοδιάστατη εκδοχή. Πρέπει να επιστρέψει την ακολουθία $P$ μήκους $M$ -- το μονοδιάστατο πρότυπο τιμών όρασης που θα επαναλαμβάνεται. Το $M$ και όλα τα στοιχεία του $P$ πρέπει να είναι μεταξύ $1$ και $60$.

\begin{lstlisting}
 int getMove1d(std::vector<int> v)
\end{lstlisting}

Η συνάρτηση αυτή θα κληθεί μία φορά για κάθε δυνατή σάρωση, αν το τεστ ελέγχου αφορά τη μονοδιάστατη εκδοχή. Δέχεται ως όρισμα μια σάρωση που παράγεται από έναν φάρο, ως ακολουθία μήκους $2V + 1$ (όπου $V$ είναι η ισχύς όρασης του φάρου στον κεντρικό τομέα της ακολουθίας). Πρέπει να επιστρέψει σε ποιον τομέα θα τηλεμεταφερθεί το σκάφος, ως δείκτη $T$ μέσα στην ακολουθία, έτσι ώστε $0 \leq T < 2V + 1$.

\begin{lstlisting}
 std::vector<std::vector<int>> getVisionPattern2d()
\end{lstlisting}

Η συνάρτηση αυτή θα κληθεί μία φορά, αν το τεστ ελέγχου αφορά τη δισδιάστατη εκδοχή. Πρέπει να επιστρέψει έναν ορθογώνιο πίνακα $P$ διαστάσεων $N$ επί $M$ -- το δισδιάστατο πρότυπο τιμών όρασης που θα επαναλαμβάνεται. Τα $N$, $M$ και όλα τα στοιχεία του $P$ πρέπει να είναι μεταξύ $1$ και $60$.

\begin{lstlisting}
 std::pair<int, int> getMove2d(std::vector<std::vector<int>> v)
\end{lstlisting}

Η συνάρτηση αυτή θα κληθεί μία φορά για κάθε δυνατή σάρωση, αν η περίπτωση ελέγχου αφορά τη δισδιάστατη εκδοχή. Δέχεται ως όρισμα μια σάρωση που παράγεται από έναν φάρο, ως τετραγωνικό πίνακα διαστάσεων $2V + 1$ επί $2V + 1$ (όπου $V$ είναι η ισχύς όρασης του φάρου στον κεντρικό τομέα του τετραγώνου). Πρέπει να επιστρέψει σε ποιον τομέα θα τηλεμεταφερθεί το σκάφος, ως ζεύγος δεικτών $S$ και $T$ μέσα στον πίνακα, έτσι ώστε $0 \leq S, T < 2V + 1$.

Για τη δοθείσα διάσταση $D$ ($1$ ή $2$) του τεστ ελέγχου, ο βαθμολογητής θα ελέγξει ότι το πρόγραμμά σας παράγει ένα έγκυρο πρότυπο -- ότι κάνει έγκυρες επιλογές τηλεμεταφοράς για όλες τις δυνατές σαρώσεις και ότι η λύση σας λειτουργεί ανεξάρτητα από τις αρχικές συντεταγμένες και από την ακολουθία προσανατολισμών των σαρώσεων.

\subsection{Περιορισμοί}

\vspace{0.1em}

\begin{itemize}
\item $1 \leq D \leq 2$
\item $1 \leq N, M, V_{X, Y}, V_Y \leq 60$
\end{itemize}

\newpage

\subsection{Υποπροβλήματα και βαθμολόγηση}

\begin{table}[H]
\begin{tblr}{|Q[c,m]|Q[c,m]|Q[c,m]|Q[c,m]|}
\hline
\textbf{Υποπρόβλημα} & \textbf{Βαθμοί} & \textbf{$D$} & \textbf{$A^*$} \\
\hline
$1$ & $24$ & $1$ & $1.5$ \\
\hline
$2$ & $76$ & $2$ & $1.125$ \\
\hline
\end{tblr}
\end{table}
\FloatBarrier

Η βαθμολογία σας για ένα συγκεκριμένο υποπρόβλημα (εφόσον η λύση σας για αυτό είναι σωστή) εξαρτάται από τη μέση ισχύ όρασης στο πρότυπο $P$ που δώσατε για αυτό. Έστω ότι αυτή είναι $A$, και ότι ο στόχος για το υποπρόβλημα είναι $A^*$. Τότε η βαθμολογία σας (ως κλάσμα των βαθμών του υποπροβλήματος) θα είναι:

$$1 - {\left(1 - \frac{A^* - 1}{\max{\left(A, A^*\right)} - 1}\right)}^{0.8}$$

\subsection{Παράδειγμα προτύπου}

Έστω ότι η λύση σας επιστρέφει το ακόλουθο πρότυπο με $N=3$ και $M=2$:

\[
\left[
\begin{array}{c c}
1 & 2 \\
3 & 4 \\
5 & 6
\end{array}
\right]
\]

Τώρα ας εξετάσουμε τις δυνατές σαρώσεις στον τομέα $(0, 1)$, ο οποίος έχει ισχύ όρασης $2$. Υπάρχουν $4$ δυνατοί προσανατολισμοί (μερικοί από τους οποίους παράγουν την ίδια σάρωση για αυτό το πρότυπο):

\[
\begin{array}{cc}
\text{Προσανατολισμός $0$: Καμία αντιστροφή} & \text{Προσανατολισμός $1$: Αντιστροφή άξονα $Y$} \\
\left[
\begin{array}{ccccc}
4 & 3 & 4 & 3 & 4 \\
6 & 5 & 6 & 5 & 6 \\
2 & 1 & 2 & 1 & 2 \\
4 & 3 & 4 & 3 & 4 \\
6 & 5 & 6 & 5 & 6
\end{array}
\right]
&
\left[
\begin{array}{ccccc}
4 & 3 & 4 & 3 & 4 \\
6 & 5 & 6 & 5 & 6 \\
2 & 1 & 2 & 1 & 2 \\
4 & 3 & 4 & 3 & 4 \\
6 & 5 & 6 & 5 & 6
\end{array}
\right]
\\
\\
\text{Προσανατολισμός $2$: Αντιστροφή άξονα $X$} & \text{Προσανατολισμός $3$: Αντιστροφή των δύο αξόνων} \\
\left[
\begin{array}{ccccc}
6 & 5 & 6 & 5 & 6 \\
4 & 3 & 4 & 3 & 4 \\
2 & 1 & 2 & 1 & 2 \\
6 & 5 & 6 & 5 & 6 \\
4 & 3 & 4 & 3 & 4
\end{array}
\right]
&
\left[
\begin{array}{ccccc}
6 & 5 & 6 & 5 & 6 \\
4 & 3 & 4 & 3 & 4 \\
2 & 1 & 2 & 1 & 2 \\
6 & 5 & 6 & 5 & 6 \\
4 & 3 & 4 & 3 & 4
\end{array}
\right]
\end{array}
\] 

Εφόσον οι προσανατολισμοί $0$ και $1$ είναι ίδιοι και το σκάφος είναι ντετερμινιστικό, θα γίνει μόνο μία κλήση της \verb|getMove2d| και για τους δύο. Έστω ότι η λύση σας επιστρέφει ως κίνηση το $(0, 1)$ για αυτήν τη σάρωση. Στον προσανατολισμό $0$ αυτό σημαίνει μετακίνηση $1$ θέσης προς τα αριστερά και $2$ προς τα πάνω, ενώ στον προσανατολισμό $1$ σημαίνει μετακίνηση $1$ θέσης προς τα δεξιά και $2$ προς τα πάνω.

\subsection{Ενδεικτικός βαθμολογητής}

Η πρώτη είσοδος του ενδεικτικού βαθμολογητή είναι το $D$. Μετά από αυτό θα εκτελέσει όλες τις κλήσεις προς τις συναρτήσεις σας: πρώτα θα πάρει το πρότυπο $P$ και έπειτα θα αποθηκεύσει όλες τις επιλογές τηλεμεταφοράς. Μετά από αυτό, θα ξεκινήσει μια αλληλεπίδραση μέσω κονσόλας, ζητώντας τις αρχικές συντεταγμένες του σκάφους. Στη συνέχεια, θα εμφανίσει την τρέχουσα κατάσταση: τις συντεταγμένες και την ``ακατέργαστη'' (χωρίς αντιστροφές) σάρωση. Έπειτα, θα ζητήσει τον προσανατολισμό που θα χρησιμοποιηθεί ($0$ για καμία αντιστροφή, $1$ για αντιστροφή του άξονα $Y$, $2$ για αντιστροφή του άξονα $X$ και $3$ για αντιστροφή και των δύο αξόνων). Μετά από αυτό, θα εμφανίσει τη σάρωση με τον δοθέντα προσανατολισμό, καθώς και την κίνηση που επιλέγει το σκάφος. Το σκάφος θα μετακινηθεί και η διαδικασία αυτή θα επαναλαμβάνεται. Παρατηρήστε ότι ο ενδεικτικός βαθμολογητής δεν θα ελέγχει αυτόματα αν η λύση σας είναι σωστή.

\end{document}